\documentclass[11pt,a4paper]{article}

\usepackage[utf8]{inputenc}
\usepackage[T1]{fontenc}
\usepackage[spanish,es-noshorthands]{babel}
\usepackage{amsmath,amssymb,amsfonts}
\usepackage[a4paper,margin=2.2cm]{geometry}
\usepackage{enumitem}
\usepackage{xcolor}
\usepackage{tcolorbox}
\tcbuselibrary{skins,breakable}
\usepackage{titlesec}
\usepackage{hyperref}
\hypersetup{colorlinks=true, linkcolor=blue!50!black, urlcolor=blue!50!black, citecolor=black,
  pdftitle={Resumen Parciales 1 y 2 - Geometria Analitica}, pdfauthor={Facundo Luna}}

\titleformat{\section}{\Large\bfseries\color{blue!35!black}}{\thesection.}{0.5em}{}
\titleformat{\subsection}{\large\bfseries}{\thesubsection.}{0.5em}{}
\titlespacing*{\section}{0pt}{1.6em}{0.8em}
\titlespacing*{\subsection}{0pt}{1.2em}{0.5em}

\newtcolorbox{teoria}[1][]{colback=blue!4!white, colframe=blue!40!black, fonttitle=\bfseries,
  breakable, title={#1}}
\newtcolorbox{proc}[1][Procedimiento]{colback=green!5!white, colframe=green!30!black, fonttitle=\bfseries,
  breakable, title={#1}}
\newtcolorbox{atencion}[1][Atención / trampas frecuentes]{colback=red!4!white, colframe=red!55!black,
  fonttitle=\bfseries, breakable, title={#1}}
\newtcolorbox{ejemplo}[1][Ejemplo resuelto]{colback=yellow!6!white, colframe=orange!55!black,
  fonttitle=\bfseries, breakable, title={#1}}

\newcommand{\qpide}[1]{\textbf{¿Qué pide el ejercicio?} #1\par}

\title{\textbf{\Huge Resumen Teórico-Práctico}\\[0.3cm]
\Large Geometría Analítica --- Parciales 1 y 2\\[0.2cm]
\normalsize Unidades I a IV: Espacios vectoriales, Transformaciones lineales,\\
Álgebra vectorial y Cónicas/Cuádricas}
\author{Facundo Luna\\ \small Universidad Católica de Salta}
\date{\today}

\begin{document}
\maketitle
\tableofcontents
\newpage

\begin{teoria}[Cómo usar este resumen]
Cada tema está organizado como un \textbf{``tipo de ejercicio''}: primero un recordatorio
teórico mínimo, después \textbf{qué pide} el enunciado, después el \textbf{procedimiento}
paso a paso para resolverlo, un \textbf{ejemplo resuelto} y, por último, una caja roja con
\textbf{errores típicos} a los que hay que prestar atención en el parcial. La idea es que,
frente a un ejercicio nuevo, puedas reconocer el ``tipo'' y aplicar directamente los pasos.
\end{teoria}

\part{PRIMER PARCIAL --- Espacios vectoriales y Transformaciones lineales}

\section{Unidad I: Espacios vectoriales}

\subsection{Determinar si un conjunto es un espacio vectorial}

\begin{teoria}[Definición]
$V$ con operaciones $+$ y $\cdot$ (escalar) es un \textbf{espacio vectorial} si se cumplen
10 axiomas: cerradura de la suma, conmutativa, asociativa, existencia de neutro $\odot$,
existencia de inverso aditivo, cerradura del producto por escalar, y las dos distributivas,
asociatividad de escalares y $1x=x$.
\end{teoria}

\qpide{Dado un conjunto $V$ con ciertas operaciones (usuales o definidas ``raro''),
decidir si es un espacio vectorial.}

\begin{proc}
\begin{enumerate}
    \item Identificar el vector nulo candidato. Si $V$ ni siquiera lo contiene, \textbf{no}
    es espacio vectorial (fin del ejercicio).
    \item Verificar cerradura de la suma: tomar $u,v\in V$ genéricos, sumarlos y ver si el
    resultado sigue perteneciendo a $V$.
    \item Verificar cerradura del producto por escalar: tomar $u\in V$, $\alpha$ escalar, y
    ver si $\alpha u\in V$.
    \item Si ambas cerraduras se cumplen, los 8 axiomas restantes casi siempre se heredan
    directamente de las propiedades de $\mathbb R$ (son rutina algebraica).
    \item Si se sospecha que \textbf{no} es espacio vectorial, conviene buscar un
    \textbf{contraejemplo numérico} concreto en vez de trabajar en general.
\end{enumerate}
\end{proc}

\begin{ejemplo}
¿Es espacio vectorial el conjunto de matrices $2\times2$ antisimétricas ($A^t=-A$)?

\textbf{Neutro:} $0^t=-0$, la matriz nula es antisimétrica $\Rightarrow$ pertenece.

\textbf{Suma:} si $A^t=-A$, $B^t=-B$: $(A+B)^t=A^t+B^t=-A-B=-(A+B)$. Cerrado.

\textbf{Escalar:} $(\alpha A)^t=\alpha A^t=\alpha(-A)=-(\alpha A)$. Cerrado.

Como ambas cerraduras valen y los demás axiomas se heredan de $M_{2\times2}$,
\textbf{sí} es espacio vectorial (de hecho, subespacio de $M_{2\times2}$).
\end{ejemplo}

\begin{atencion}
\begin{itemize}
    \item Un solo axioma que falle alcanza para descartar todo el conjunto.
    \item El conjunto de polinomios de grado \textbf{exactamente} $n$ (no ``menor o igual'')
    \textbf{no} es espacio vectorial: la suma de dos de grado $n$ puede bajar de grado
    (cancelación del término principal), y el polinomio nulo (que debe estar) no tiene
    grado $n$. Es un clásico ``truco'' de examen.
    \item Operaciones ``raras'' del tipo $\alpha(x,y)=(\alpha x,0)$ suelen fallar en el
    axioma $1\cdot x=x$, no en la cerradura: conviene probar ese axioma primero si la
    multiplicación por escalar se ve sospechosa.
\end{itemize}
\end{atencion}

\subsection{Determinar si un subconjunto es subespacio}

\begin{teoria}[Teorema 1.1.4]
$H\subseteq V$ no vacío es subespacio si y solo si:
\textbf{(i)} $x,y\in H\Rightarrow x+y\in H$ y \textbf{(ii)} $x\in H\Rightarrow \alpha x\in H$
$\forall\alpha$. Todo subespacio contiene al $\odot$.
\end{teoria}

\qpide{Dado $H$ (subconjunto de un espacio conocido, definido por una ecuación o
condición), decidir si es subespacio.}

\begin{proc}
\begin{enumerate}
    \item \textbf{Atajo:} verificar primero si $0\in H$. Si no, \textbf{no} es subespacio
    (se termina el ejercicio).
    \item Si $0\in H$: tomar $u,v$ genéricos de $H$ y probar $u+v\in H$.
    \item Probar $\alpha u\in H$ para $\alpha$ escalar genérico.
    \item Si ambas valen, $H$ es subespacio. Si se sospecha que no, buscar un
    contraejemplo numérico simple (probar con valores concretos).
\end{enumerate}
\end{proc}

\begin{ejemplo}
$H=\{(x,y,z)\in\mathbb R^3: x-2y+z=0\}$.

$0$: $(0,0,0)$ cumple $0-0+0=0$ \checkmark.

Suma: $(x_1-2y_1+z_1)+(x_2-2y_2+z_2)=0+0=0\Rightarrow$ la suma está en $H$.

Escalar: $\alpha(x-2y+z)=\alpha\cdot0=0\Rightarrow$ está en $H$.

Es subespacio. Como es 1 ecuación homogénea con 3 incógnitas, $\dim H=3-1=2$. Despejando
$y=\frac{x+z}{2}$: $(x,\frac{x+z}2,z)=x(1,\tfrac12,0)+z(0,\tfrac12,1)$, con base (sin
fracciones) $\{(2,1,0),(0,1,2)\}$.
\end{ejemplo}

\begin{atencion}
\begin{itemize}
    \item Ecuaciones con término independiente $\neq0$ (ej. $x+y+z=1$) \textbf{nunca} dan
    subespacio: el $(0,0,0)$ no las cumple. Revisar esto primero ahorra tiempo.
    \item Condiciones con valor absoluto o cuadrados ($|y|=|x|$, $y=x^2$) suelen fallar en
    la cerradura aunque contengan al $0$: buscar un contraejemplo numérico chico.
    \item El conjunto de matrices con $\det=0$ contiene al $0$ pero \textbf{falla en la
    suma} (suma de dos singulares puede no ser singular): contraejemplo clásico a
    recordar de memoria.
\end{itemize}
\end{atencion}

\subsection{Combinación lineal y espacio generado}

\qpide{(a) Decidir si un vector $w$ es combinación lineal de otros dados (equivalente a
$w\in\operatorname{gen}\{v_1,\dots,v_k\}$). (b) Hallar la ecuación cartesiana del
subespacio generado por un conjunto de vectores.}

\begin{proc}[Procedimiento (a)]
\begin{enumerate}
    \item Plantear $w=\alpha_1v_1+\dots+\alpha_kv_k$.
    \item Igualar componente a componente $\Rightarrow$ sistema lineal en $\alpha_i$.
    \item Resolver: si es \textbf{compatible}, $w$ es combinación lineal (dar los
    $\alpha_i$); si es \textbf{incompatible}, no lo es.
\end{enumerate}
\end{proc}

\begin{proc}[Procedimiento (b), cuando hay más componentes que vectores generadores]
\begin{enumerate}
    \item Plantear $(x,y,z,\dots)=\alpha_1v_1+\dots+\alpha_kv_k$.
    \item Armar la matriz aumentada con $(x,y,z,\dots)$ como términos independientes y
    escalonar por Gauss.
    \item Al quedar más ecuaciones que incógnitas $\alpha$, las filas que se anulan del
    lado de los $\alpha$ dan, igualadas a $0$, la(s) ecuación(es) cartesiana(s) del
    espacio generado.
\end{enumerate}
\end{proc}

\begin{atencion}
\begin{itemize}
    \item Si la cantidad de vectores generadores coincide con la dimensión del espacio
    ambiente (ej. 3 vectores en $\mathbb R^3$) y son LI, el generado es \textbf{todo} el
    espacio (no hay ecuación restrictiva): alcanza con un determinante, no hace falta el
    sistema.
    \item No confundir ``generar'' (alcanza con que todo vector se escriba como
    combinación lineal) con ``ser base'' (además deben ser LI).
\end{itemize}
\end{atencion}

\subsection{Dependencia e independencia lineal}

\begin{teoria}[Definición]
$v_1,\dots,v_n$ son LI si $\alpha_1v_1+\dots+\alpha_nv_n=0$ solo se cumple con todos los
$\alpha_i=0$. Si existe una combinación no trivial que da $0$, son LD.
\end{teoria}

\qpide{Determinar si un conjunto de vectores (o polinomios, o matrices) es LI o LD.}

\begin{proc}
\begin{enumerate}
    \item Plantear $\alpha_1v_1+\dots+\alpha_nv_n=0$ y armar el sistema homogéneo.
    \item \textbf{Atajo con determinante:} si son $n$ vectores en $\mathbb R^n$, armar la
    matriz $A$ con los vectores como columnas: $\det A\neq0\Rightarrow$ LI;
    $\det A=0\Rightarrow$ LD.
    \item \textbf{Caso general} (matrices, polinomios, o cantidad $\neq$ dimensión):
    escalonar por Gauss. Solución única (trivial) $\Rightarrow$ LI; infinitas soluciones
    $\Rightarrow$ LD (exhibir una combinación no trivial dando un valor libre, p.ej.\ $1$).
\end{enumerate}
\end{proc}

\begin{ejemplo}
¿$(1,-2,3),(2,-2,0),(0,1,7)$ son LI? Se arma $\alpha(1,-2,3)+\beta(2,-2,0)+\gamma(0,1,7)=0$,
que da el sistema $\alpha+2\beta=0$, $-2\alpha-2\beta+\gamma=0$, $3\alpha+7\gamma=0$.
Escalonando se llega a $\alpha=\beta=\gamma=0$ (única solución) $\Rightarrow$ \textbf{LI}.
\end{ejemplo}

\begin{atencion}
\begin{itemize}
    \item \textbf{Teorema 1.2.3:} un conjunto de $n$ vectores en $\mathbb R^m$ con $n>m$
    es \textbf{siempre} LD (más vectores que ``grados de libertad''). Revisar esto primero
    puede evitar todo el cálculo.
    \item Con solo 2 vectores, alcanza con ver si uno es múltiplo escalar del otro.
    \item Solo importa si el determinante da \textbf{cero o no}, no su valor ni su signo.
\end{itemize}
\end{atencion}

\subsection{Base y dimensión}

\begin{teoria}[Teoremas clave]
\textbf{1.2.15:} Cualquier conjunto de $n$ vectores LI en un espacio de dimensión $n$ es
automáticamente una base (no hace falta probar que genera). \textbf{1.2.13:} en un espacio
de dimensión $n$, no puede haber más de $n$ vectores LI.
\end{teoria}

\qpide{(a) Verificar si un conjunto dado es base. (b) Hallar una base y la dimensión de un
subespacio dado por ecuación(es).}

\begin{proc}[Procedimiento (a)]
Si el conjunto tiene exactamente $\dim V=n$ vectores: \textbf{alcanza con probar que son
LI} (determinante $\neq0$). Si tiene menos que $n$: nunca es base (no alcanza a generar).
Si tiene más que $n$: nunca es base (no pueden ser LI).
\end{proc}

\begin{proc}[Procedimiento (b)]
\begin{enumerate}
    \item Despejar las variables ``dependientes'' en función de las ``libres'' usando la(s)
    ecuación(es).
    \item Escribir el vector genérico agrupando por cada variable libre (factor común).
    \item Cada vector que multiplica a una variable libre es un generador; juntos forman
    una base (son LI automáticamente).
    \item $\dim=$ cantidad de variables libres $= n - (\text{ecuaciones independientes})$.
\end{enumerate}
\end{proc}

\begin{atencion}
\begin{itemize}
    \item Con más de una ecuación, verificar que sean independientes entre sí antes de
    contar los grados de libertad.
    \item Una base \textbf{nunca} es única: si tu resultado no coincide ``literalmente''
    con el de otra persona, verificá que genere el mismo subespacio antes de asumir un
    error.
\end{itemize}
\end{atencion}

\subsection{Coordenadas y matriz de transición (cambio de base)}

\qpide{Dado un vector en una base, hallar sus coordenadas en otra base; o construir la
matriz de cambio de base entre $\beta_1$ y $\beta_2$.}

\begin{proc}[Caso: de la base canónica $\beta_1$ a otra base $\beta_2=\{v_1,\dots,v_n\}$]
\begin{enumerate}
    \item Armar $C$ con los $v_i$ (en coordenadas canónicas) como columnas. $C$ es la
    matriz de transición \textbf{de $\beta_2$ a $\beta_1$}.
    \item $A=C^{-1}$ es la matriz de transición \textbf{de $\beta_1$ a $\beta_2$}:
    $[x]_{\beta_2}=A\,[x]_{\beta_1}$.
\end{enumerate}
\end{proc}

\begin{proc}[Caso general: ninguna de las dos es canónica]
\begin{enumerate}
    \item Expresar cada vector de $\beta_1$ como combinación lineal de los de $\beta_2$
    (un sistema por cada vector de $\beta_1$).
    \item Los coeficientes hallados para cada vector de $\beta_1$ forman una columna de la
    matriz de transición $A$ de $\beta_1$ a $\beta_2$.
    \item $[x]_{\beta_2}=A\,[x]_{\beta_1}$.
\end{enumerate}
\end{proc}

\begin{atencion}
\begin{itemize}
    \item Es muy fácil confundir el sentido: la matriz de transición \textbf{de $\beta_1$
    a $\beta_2$} se arma expresando los vectores \textbf{de $\beta_1$} \textbf{en
    términos de} $\beta_2$.
    \item \textbf{Teorema 1.3.4:} si $A$ es la transición de $\beta_1$ a $\beta_2$,
    $A^{-1}$ es la de $\beta_2$ a $\beta_1$ — evita rehacer cuentas.
    \item Para verificar independencia lineal en espacios abstractos (polinomios,
    matrices), convertir todo a coordenadas respecto de una base conocida y aplicar el
    criterio del determinante (Teorema 1.3.5).
\end{itemize}
\end{atencion}

\newpage
\section{Unidad II: Transformaciones lineales}

\subsection{Determinar si una transformación es lineal}

\begin{teoria}[Definición]
$T:V\to W$ es lineal si $T(u+v)=Tu+Tv$ y $T(\alpha v)=\alpha Tv$ para todo $u,v,\alpha$.
\end{teoria}

\qpide{Dada la fórmula de $T$, decidir si es una transformación lineal.}

\begin{proc}
\begin{enumerate}
    \item \textbf{Atajo rápido:} calcular $T(0)$. Si $T(0)\neq0$ (el $0$ de $W$),
    \textbf{no} es lineal (fin del ejercicio, sin más cuentas).
    \item Si $T(0)=0$: verificar aditividad, $T(u+v)\stackrel{?}{=}T(u)+T(v)$.
    \item Verificar homogeneidad, $T(\alpha u)\stackrel{?}{=}\alpha T(u)$.
    \item Si ambas valen, $T$ es lineal.
\end{enumerate}
\end{proc}

\begin{ejemplo}
$T:\mathbb R^2\to\mathbb R^{2\times2}$, $T(x,y)=\begin{pmatrix}x-1&y\\2x&x+y\end{pmatrix}$.

$T(0,0)=\begin{pmatrix}-1&0\\0&0\end{pmatrix}\neq\begin{pmatrix}0&0\\0&0\end{pmatrix}$
$\Rightarrow$ \textbf{no es lineal} (no hace falta seguir).
\end{ejemplo}

\begin{atencion}
\begin{itemize}
    \item Cualquier término \textbf{constante}, \textbf{cuadrático} ($x^2$), \textbf{producto
    cruzado} ($xy$) o \textbf{valor absoluto/trigonométrico} de las variables casi siempre
    rompe la linealidad. Mirar la fórmula antes de hacer cuentas: si hay un término
    independiente, alcanza con evaluar $T(0)$.
    \item $f(x)=mx+b$ es una función ``lineal'' en el sentido de recta, pero \textbf{no}
    es transformación lineal si $b\neq0$.
\end{itemize}
\end{atencion}

\subsection{Núcleo, imagen, dimensión y clasificación}

\begin{teoria}[Definiciones y teorema de la dimensión]
$\operatorname{Nu}(T)=\{v: Tv=0\}$, $\operatorname{Im}(T)=\{Tv: v\in V\}$. Ambos son
subespacios. \textbf{Teorema:} $\dim V=\dim\operatorname{Nu}(T)+\dim\operatorname{Im}(T)$.
$T$ inyectiva (monomorfismo) $\Leftrightarrow \operatorname{Nu}(T)=\{0\}$. $T$ sobreyectiva
(epimorfismo) $\Leftrightarrow \dim\operatorname{Im}(T)=\dim W$. Isomorfismo = ambas.
\end{teoria}

\qpide{Hallar $\operatorname{Nu}(T)$ (base y dimensión), $\dim\operatorname{Im}(T)$, y
clasificar $T$.}

\begin{proc}
\begin{enumerate}
    \item Plantear $T(v)=0$ (el $0$ de $W$) y resolver el sistema en las variables del
    dominio.
    \item Agrupar por variables libres para obtener una base de $\operatorname{Nu}(T)$
    (igual técnica que en Base y dimensión); $\dim\operatorname{Nu}(T)=$ cantidad de
    libres.
    \item Teorema de la dimensión $\Rightarrow \dim\operatorname{Im}(T)=\dim V-\dim
    \operatorname{Nu}(T)$.
    \item Clasificar: $\operatorname{Nu}(T)=\{0\}\Rightarrow$ mono;
    $\dim\operatorname{Im}(T)=\dim W\Rightarrow$ epi; ambas $\Rightarrow$ iso.
\end{enumerate}
\end{proc}

\begin{atencion}
\begin{itemize}
    \item \textbf{Teorema 2.2.4:} si $\dim V>\dim W$, $T$ \textbf{nunca} es inyectiva; si
    $\dim V<\dim W$, \textbf{nunca} es sobreyectiva. Revisar dimensiones antes de calcular
    ahorra trabajo.
    \item \textbf{Teorema 2.2.3:} si $\dim V=\dim W$, inyectiva $\Leftrightarrow$
    sobreyectiva — probar una de las dos alcanza.
    \item Se pide una \textbf{base} del núcleo, no una ecuación: dar vectores generadores.
\end{itemize}
\end{atencion}

\subsection{Expresión general de un morfismo dado por sus valores}

\qpide{Dados $T(v_1),\dots,T(v_n)$ para una base $\{v_1,\dots,v_n\}$ del dominio, hallar la
fórmula $T(x,y,z,\dots)$.}

\begin{proc}
\begin{enumerate}
    \item Verificar que $\{v_1,\dots,v_n\}$ es base del dominio (determinante $\neq0$).
    \item Escribir $(x,y,\dots)=\alpha_1v_1+\dots+\alpha_nv_n$ y resolver para los
    $\alpha_i$ en función de $x,y,\dots$
    \item Por linealidad: $T(x,y,\dots)=\alpha_1T(v_1)+\dots+\alpha_nT(v_n)$.
    \item \textbf{Verificar} sustituyendo los datos originales (debe reproducirlos).
\end{enumerate}
\end{proc}

\begin{atencion}
\begin{itemize}
    \item Si el conjunto dado \textbf{no} es base del dominio, $T$ no queda unívocamente
    determinada — revisar el enunciado.
    \item El paso de verificación final es la mejor forma de pescar un error de cuenta
    antes de seguir con la clasificación.
\end{itemize}
\end{atencion}

\subsection{Matriz asociada a una transformación lineal}

\begin{teoria}
$A_T[v]_{\beta_1}=[Tv]_{\beta_2}$, donde las columnas de $A_T$ son
$[Tv_1]_{\beta_2},\dots,[Tv_n]_{\beta_2}$ con $\{v_1,\dots,v_n\}=\beta_1$.
\end{teoria}

\qpide{Hallar la matriz de $T$ respecto de bases dadas y usarla para hallar la imagen de
un vector.}

\begin{proc}[Bases canónicas]
Aplicar $T$ a cada vector de la base canónica del dominio; cada resultado (ya en
coordenadas canónicas del codominio) es una columna de $[T]$.
\end{proc}

\begin{proc}[Bases arbitrarias $\beta_1$ (dominio), $\beta_2$ (codominio)]
\begin{enumerate}
    \item Aplicar $T$ a cada vector de $\beta_1$.
    \item Expresar cada resultado como combinación lineal de $\beta_2$ (un sistema por
    cada uno).
    \item Esos coeficientes forman las columnas de $[T]_{\beta_1}^{\beta_2}$.
    \item Para la imagen de $v$ (dado en coordenadas estándar): hallar $[v]_{\beta_1}$,
    multiplicar $[T]_{\beta_1}^{\beta_2}\cdot[v]_{\beta_1}=[Tv]_{\beta_2}$, y si hace falta,
    reconstruir en coordenadas estándar combinando los vectores de $\beta_2$ con esos
    coeficientes.
\end{enumerate}
\end{proc}

\begin{atencion}
\begin{itemize}
    \item El orden de las columnas sigue el orden de la base del dominio: si se cambia el
    orden de la base, cambia el orden de las columnas.
    \item Siempre conviene \textbf{verificar} calculando $T(v)$ directamente por la
    fórmula y comparando con el resultado matricial.
    \item No confundir $[Tv]_{\beta_2}$ (coordenadas en $\beta_2$) con $Tv$ en coordenadas
    estándar: son representaciones distintas del mismo vector.
\end{itemize}
\end{atencion}

\subsection{Autovalores y autovectores}

\begin{teoria}[Definiciones y teoremas]
$Av=\lambda v$ con $v\neq0$. Ecuación característica: $p(\lambda)=\det(A-\lambda I)=0$. La
multiplicidad geométrica ($\dim E_\lambda$) es siempre $\le$ la algebraica. $A$ es
diagonalizable $\Leftrightarrow$ para cada $\lambda$, mult.\ geométrica = mult.\ algebraica
(en particular, si todos los autovalores son distintos, siempre es diagonalizable).
\end{teoria}

\qpide{Hallar polinomio característico, autovalores, autovectores/autoespacios y decidir
diagonalizabilidad.}

\begin{proc}
\begin{enumerate}
    \item Plantear $p(\lambda)=\det(A-\lambda I)=0$.
    \item Resolver para hallar $\lambda_1,\dots,\lambda_k$ (con sus multiplicidades
    algebraicas).
    \item Para cada $\lambda_i$: resolver $(A-\lambda_iI)v=0$; las soluciones no
    triviales forman $E_{\lambda_i}$, su dimensión es la mult.\ geométrica.
    \item Diagonalizable $\Leftrightarrow$ mult.\ geométrica = mult.\ algebraica para
    \textbf{todos} los autovalores.
\end{enumerate}
\end{proc}

\begin{ejemplo}
$B=\begin{pmatrix}2&1\\6&3\end{pmatrix}$. $\det(B-\lambda I)=(2-\lambda)(3-\lambda)-6=
\lambda^2-5\lambda=\lambda(\lambda-5)$. Autovalores $\lambda_1=0,\lambda_2=5$ (distintos)
$\Rightarrow$ diagonalizable automáticamente. Para $\lambda_1=0$: $2x+y=0\Rightarrow
v_1=(1,-2)$. Para $\lambda_2=5$: $-3x+y=0\Rightarrow v_2=(1,3)$.
\end{ejemplo}

\begin{atencion}
\begin{itemize}
    \item Si $A$ es triangular (o diagonal), los autovalores son directamente los
    elementos de la diagonal — no hace falta el determinante.
    \item Al resolver $(A-\lambda I)v=0$ siempre debe aparecer \textbf{al menos una fila de
    ceros} al escalonar; si no aparece, hay un error en el polinomio característico o en
    la resolución.
    \item Si para un autovalor de mult.\ algebraica 2 el autoespacio da dimensión 1, la
    matriz \textbf{no} es diagonalizable — no hace falta seguir con los demás autovalores.
    \item \textbf{Verificación rápida:} suma de autovalores = traza($A$); producto de
    autovalores = $\det(A)$. Chequear esto antes de seguir ahorra errores.
\end{itemize}
\end{atencion}

\subsection{Diagonalización y matrices semejantes}

\begin{teoria}
$A$ diagonalizable si tiene $n$ autovectores LI: $D=P^{-1}AP$, con $P$ formada por los
autovectores como columnas y $D$ diagonal con los autovalores en el mismo orden.
\end{teoria}

\qpide{Hallar $P$ y $D$, y verificar la semejanza.}

\begin{proc}
\begin{enumerate}
    \item Obtener autovalores y una base de cada autoespacio.
    \item Armar $P$ con los autovectores como columnas (en cualquier orden, pero fijo).
    \item Armar $D$ con los autovalores en la diagonal, \textbf{en el mismo orden} que las
    columnas de $P$.
    \item Para verificar sin invertir $P$: comprobar $AP=PD$ columna a columna (equivale a
    $Av_i=\lambda_iv_i$, ya verificado al calcular los autovectores).
\end{enumerate}
\end{proc}

\begin{atencion}
El orden de columnas de $P$ y de la diagonal de $D$ \textbf{deben coincidir} — es el error
más común de este tipo de ejercicio. No hace falta calcular $P^{-1}$ para argumentar la
semejanza: basta mostrar $AP=PD$ y que $P$ es invertible (columnas LI, provenientes de
autoespacios distintos).
\end{atencion}

\newpage
\part{SEGUNDO PARCIAL --- Álgebra vectorial, rectas/planos y Cónicas/Cuádricas}

\section{Unidad III: Álgebra vectorial y geometría del espacio}

\subsection{Producto escalar, proyección, producto vectorial y producto mixto}

\begin{teoria}[Fórmulas]
\begin{align*}
\text{Proyección: } &\operatorname{proy}_v u=\frac{u\cdot v}{|v|^2}\,v\\
\text{Producto vectorial: } &u\times v=\begin{vmatrix}i&j&k\\a_1&b_1&c_1\\a_2&b_2&c_2\end{vmatrix},
\qquad |u\times v|=|u||v|\sen\varphi\\
\text{Producto mixto: } &(u,v,w)=(u\times v)\cdot w=\det\begin{pmatrix}u\\v\\w\end{pmatrix}
\end{align*}
\end{teoria}

\qpide{Calcular proyecciones, productos vectoriales/mixtos, o usarlos para áreas,
volúmenes y ortogonalidad.}

\begin{proc}
\begin{itemize}
    \item $u\times v$ es \textbf{ortogonal} a $u$ y a $v$ (útil para hallar normales).
    \item Área del paralelogramo de lados $u,v$: $|u\times v|$. Área del
    \textbf{triángulo} $ABC$: $\frac12\left|\overrightarrow{AB}\times\overrightarrow{AC}\right|$ (¡la mitad!).
    \item Volumen del paralelepípedo de aristas $u,v,w$: $|(u,v,w)|$.
    \item Coplanaridad de 3 vectores (o 4 puntos): $(u,v,w)=0$.
\end{itemize}
\end{proc}

\begin{atencion}
\begin{itemize}
    \item El producto vectorial \textbf{solo} existe en $\mathbb R^3$ (nunca en el plano).
    \item $u\times v=-(v\times u)$: el orden importa si se pide una orientación específica
    (regla de la mano derecha).
    \item Olvidar el factor $\frac12$ en el área de un triángulo es el error más común.
    \item La proyección tiene la misma dirección que $v$ si $u\cdot v>0$, opuesta si
    $u\cdot v<0$, y es nula si $u\perp v$.
\end{itemize}
\end{atencion}

\subsection{Bases ortonormales --- Gram-Schmidt}

\qpide{A partir de una base $\{v_1,\dots,v_k\}$, construir una base ortonormal
$\{u_1,\dots,u_k\}$.}

\begin{proc}
\begin{enumerate}
    \item $u_1=v_1/|v_1|$.
    \item Para cada $v_j$ siguiente: restar la proyección sobre \textbf{todos} los $u$ ya
    calculados: $v_j'=v_j-(v_j\cdot u_1)u_1-(v_j\cdot u_2)u_2-\dots$
    \item Normalizar: $u_j=v_j'/|v_j'|$.
    \item Repetir hasta agotar los vectores.
\end{enumerate}
\end{proc}

\begin{atencion}
\begin{itemize}
    \item Al construir el tercer vector en adelante, es fácil olvidar restar la
    proyección sobre \textbf{todos} los $u$ anteriores (no solo el último).
    \item Verificación rápida: cada $u_i$ debe tener módulo 1, y el producto escalar entre
    cualquier par debe dar 0.
\end{itemize}
\end{atencion}

\subsection{Ecuaciones de la recta}

\begin{teoria}[Formas]
Vectorial: $(x,y,z)=P+tv$. Paramétrica: $x=x_0+at,\,y=y_0+bt,\,z=z_0+ct$. Simétrica (si
$a,b,c\neq0$): $\frac{x-x_0}{a}=\frac{y-y_0}{b}=\frac{z-z_0}{c}$.
\end{teoria}

\qpide{Hallar las tres formas de la ecuación de una recta a partir de puntos, un vector
director, o condiciones de perpendicularidad.}

\begin{proc}
\begin{enumerate}
    \item Vector director: si dan 2 puntos $P,Q$: $v=Q-P$. Si piden perpendicular a uno o
    dos vectores dados: $v$ = producto vectorial de esos vectores.
    \item Escribir vectorial, luego separar por componentes (paramétrica), luego despejar
    $t$ e igualar (simétrica).
\end{enumerate}
\end{proc}

\begin{atencion}
\begin{itemize}
    \item Si una componente de $v$ es 0 (p.ej.\ $b=0$), \textbf{no} se puede escribir esa
    fracción en la simétrica (división por 0): esa variable queda como ecuación aparte
    ($y=y_0$), y la simétrica se escribe solo con las otras dos componentes.
    \item Las ecuaciones de una recta \textbf{no son únicas} (dependen del punto y de la
    escala del vector director).
    \item Para verificar si un punto pertenece a una recta, el \textbf{mismo} valor de $t$
    debe satisfacer las tres ecuaciones a la vez (no alcanza con una sola componente).
\end{itemize}
\end{atencion}

\subsection{Ecuación del plano}

\begin{teoria}
Con punto $P_0$ y normal $n=(a,b,c)$: $a(x-x_0)+b(y-y_0)+c(z-z_0)=0$.
\end{teoria}

\qpide{Hallar la ecuación de un plano a partir de (a) punto+normal, (b) tres puntos, (c)
condición de paralelismo/perpendicularidad con una recta.}

\begin{proc}
\begin{itemize}
    \item \textbf{(a)} Aplicar la fórmula directamente.
    \item \textbf{(b) Tres puntos $P,Q,R$:} $n=\overrightarrow{PQ}\times\overrightarrow{PR}$,
    luego aplicar (a) con cualquiera de los tres puntos.
    \item \textbf{(c) Perpendicular a una recta dada:} el vector director de la recta
    \textbf{es} la normal del plano.
    \item \textbf{(d) Paralelo a una recta dada:} el vector director de la recta debe ser
    ortogonal a la normal ($v\cdot n=0$); hace falta un dato más para determinar el plano
    (otro punto, otra dirección, etc.).
\end{itemize}
\end{proc}

\begin{atencion}
\begin{itemize}
    \item Confundir ``perpendicular a la recta'' ($v=n$) con ``paralelo a la recta''
    ($v\cdot n=0$) es el error más frecuente de este tema.
    \item Con tres puntos, verificar que no sean colineales (si lo son,
    $\overrightarrow{PQ}\times\overrightarrow{PR}=0$ y no queda determinado un único
    plano).
    \item La ecuación de un plano no es única: dos ecuaciones representan el mismo plano
    si sus coeficientes $(a,b,c,d)$ son proporcionales.
\end{itemize}
\end{atencion}

\subsection{Posiciones relativas}

\begin{proc}[Recta $r$ (director $v_1$) y recta $s$ (director $v_2$)]
\begin{enumerate}
    \item ¿$v_1,v_2$ proporcionales?
    \begin{itemize}
        \item \textbf{Sí:} tomar un punto de $r$ y ver si está en $s$: si sí, son
        \textbf{coincidentes}; si no, \textbf{paralelas}.
        \item \textbf{No:} armar el sistema igualando las paramétricas (con parámetros
        \emph{distintos}, $t$ y $s$). Solución \textbf{única} $\Rightarrow$
        \textbf{secantes} (calcular el punto). \textbf{Incompatible} $\Rightarrow$
        \textbf{alabeadas}.
    \end{itemize}
\end{enumerate}
\end{proc}

\begin{proc}[Plano $\pi_1$ (normal $n_1$) y plano $\pi_2$ (normal $n_2$)]
¿$n_1,n_2$ proporcionales? \textbf{Sí:} punto de $\pi_1$ en $\pi_2$ $\Rightarrow$
coincidentes; si no, paralelos. \textbf{No:} secantes en una recta (resolver el sistema
de las dos ecuaciones, queda 1 variable libre).
\end{proc}

\begin{proc}[Recta $r$ (director $v$) y plano $\pi$ (normal $n$)]
Calcular $v\cdot n$. \textbf{$=0$:} paralela o contenida — probar un punto de $r$ en
$\pi$. \textbf{$\neq0$:} secantes — sustituir la paramétrica de $r$ en $\pi$, despejar
$t$, hallar el punto.
\end{proc}

\begin{atencion}
\begin{itemize}
    \item Recta-recta usa \textbf{vectores directores}; plano-plano usa \textbf{normales};
    recta-plano compara el director de la recta con la \textbf{normal} del plano (y aquí
    ``paralelo'' significa $v\perp n$, ¡al revés que en recta-recta!).
    \item ``Alabeadas'' no tiene análogo en el plano $\mathbb R^2$: solo existe en el
    espacio.
    \item Al intersecar dos rectas, usar letras de parámetro \textbf{distintas} para cada
    una; igualar con el mismo parámetro es un error muy común.
\end{itemize}
\end{atencion}

\subsection{Distancias}

\begin{teoria}[Fórmulas]
\begin{align*}
d(P,Q)&=\sqrt{(x_2-x_1)^2+(y_2-y_1)^2+(z_2-z_1)^2}\\
d(P,r)&=\frac{|\overrightarrow{PQ}\times v|}{|v|}\quad(Q\in r)\\
d(P,\pi)&=\frac{|\overrightarrow{QP}\cdot n|}{|n|}\quad(Q\in\pi)\\
d(r,s)_{\text{alabeadas}}&=\frac{|\overrightarrow{PQ}\cdot(v_1\times v_2)|}{|v_1\times v_2|}
\quad(P\in r,\,Q\in s)
\end{align*}
Rectas/planos paralelos: se reducen a distancia punto-recta / punto-plano.
\end{teoria}

\begin{atencion}
\begin{itemize}
    \item \textbf{Antes} de aplicar cualquier fórmula, verificar la posición relativa: si
    se cortan o coinciden, la distancia es $0$ directamente.
    \item Para rectas alabeadas, \textbf{nunca} usar la fórmula punto-recta: hace falta el
    producto mixto con ambos directores.
    \item El punto $Q$ elegido (de la recta o el plano) puede ser cualquiera: el resultado
    no depende de cuál se use.
\end{itemize}
\end{atencion}

\subsection{Ángulos}

\begin{teoria}[Fórmulas]
Entre rectas ($u,v$ directores) o entre planos ($n_1,n_2$ normales):
$\cos\alpha=\dfrac{|u\cdot v|}{|u||v|}$. Entre recta ($v$) y plano ($n$): primero
$\cos\beta=\dfrac{|v\cdot n|}{|v||n|}$, y el ángulo buscado es
$\boxed{\alpha=90^\circ-\beta}$.
\end{teoria}

\begin{atencion}
\begin{itemize}
    \item El ángulo recta-plano usa el \textbf{complementario} del ángulo con la normal —
    es el error más frecuente del tema. Motivo: la normal es perpendicular al plano, así
    que hay que ``compensar'' los $90^\circ$.
    \item Usar \textbf{siempre} valor absoluto en el numerador: los ángulos entre
    rectas/planos se definen como el menor posible (agudo).
    \item $\cos=0\Rightarrow$ perpendiculares; $\cos=1\Rightarrow$ paralelos/coincidentes.
\end{itemize}
\end{atencion}

\newpage
\section{Unidad IV: Cónicas y cuádricas}

\subsection{Identificación rápida de una cónica}

\begin{teoria}[Regla práctica]
Dada $Ax^2+Cy^2+Dx+Ey+F=0$ (sin término $Bxy$):
\begin{itemize}
    \item $AC>0$, $A\neq C$ $\to$ \textbf{Elipse}. \quad $A=C\neq0$ $\to$
    \textbf{Circunferencia}.
    \item $AC<0$ $\to$ \textbf{Hipérbola}.
    \item $AC=0$ (uno de los dos es 0) $\to$ \textbf{Parábola}.
\end{itemize}
\end{teoria}

\begin{proc}[Método general: completar cuadrados]
\begin{enumerate}
    \item Agrupar los términos en $x$ por un lado y en $y$ por otro.
    \item \textbf{Sacar como factor común} el coeficiente del término cuadrático de cada
    grupo \textbf{antes} de completar el cuadrado.
    \item Completar cuadrados sumando y restando $(\text{mitad del coef.\ lineal})^2$
    \textbf{dentro} del paréntesis (recordando multiplicar por el factor común sacado al
    pasar el ``resto'' al otro lado).
    \item Reordenar para dejar la ecuación en su forma canónica.
\end{enumerate}
\end{proc}

\begin{atencion}
Olvidar sacar el factor común \textbf{antes} de completar el cuadrado es el error de
cuentas más común de todo el tema (ejemplo correcto: $4x^2-8x=4(x^2-2x)=4(x^2-2x+1-1)=
4(x-1)^2-4$, no ``$4x^2-8x+1-1$'').
\end{atencion}

\subsection{Circunferencia}

\begin{teoria}
$(x-h)^2+(y-k)^2=r^2$, centro $(h,k)$, radio $r$. Con ecuación general
$x^2+y^2+Dx+Ey+F=0$: existe si $D^2+E^2-4F>0$, con centro $\left(-\frac D2,-\frac
E2\right)$ y $r=\frac{\sqrt{D^2+E^2-4F}}{2}$.
\end{teoria}

\qpide{Hallar centro y radio a partir de la ecuación general, o la ecuación a partir de
centro/radio/condiciones.}

\begin{proc}
Completar cuadrados en $x$ y en $y$ por separado (dividiendo primero por el coeficiente
de $x^2$ si no es 1) y leer directamente centro y radio.
\end{proc}

\begin{atencion}
Si tras completar cuadrados el término independiente da negativo, no hay circunferencia
real; si da 0, se reduce a un solo punto.
\end{atencion}

\subsection{Elipse}

\begin{teoria}
$\dfrac{x^2}{a^2}+\dfrac{y^2}{b^2}=1$ con $a^2=b^2+c^2$ ($a$ = semieje \textbf{mayor}, es
la ``hipotenusa''). El eje mayor (focos y vértices principales) está siempre bajo el
denominador \textbf{más grande}, sea $x$ o $y$.
\end{teoria}

\begin{proc}
\begin{enumerate}
    \item Llevar a forma canónica (dividir o completar cuadrados si el centro no es el
    origen).
    \item El denominador mayor es $a^2$; el menor, $b^2$. Si $a^2$ está bajo $x$: elipse
    horizontal; si bajo $y$: vertical.
    \item $c=\sqrt{a^2-b^2}$. Vértices a distancia $a$ del centro (eje mayor); extremos
    del eje menor a distancia $b$; focos a distancia $c$ (mismo eje que los vértices).
\end{enumerate}
\end{proc}

\begin{atencion}
El eje mayor sigue al denominador \textbf{más grande}, independientemente de si acompaña
a $x$ o a $y$ — asumir que ``$x$ siempre es horizontal=mayor'' es un error común. Si el
centro no es el origen, todo se mide desde el centro $(h,k)$, no desde el origen.
\end{atencion}

\subsection{Parábola}

\begin{teoria}
Vértice $(0,0)$: $x^2=4py$ (vertical) o $y^2=4px$ (horizontal). Vértice $(h,k)$:
$(x-h)^2=4p(y-k)$ o $(y-k)^2=4p(x-h)$. Distancia vértice-foco = vértice-directriz = $|p|$.
\end{teoria}

\begin{proc}
\begin{enumerate}
    \item Completar cuadrados \textbf{solo en la variable que aparece al cuadrado} (la
    otra queda lineal).
    \item Identificar $4p$ (coeficiente de la variable lineal) y despejar $p$.
    \item Orientación: según qué variable está al cuadrado y el signo de $p$ (positivo:
    hacia arriba/derecha; negativo: hacia abajo/izquierda).
\end{enumerate}
\end{proc}

\begin{atencion}
Solo \textbf{una} variable aparece al cuadrado (a diferencia de elipse e hipérbola). Si al
completar cuadrados quedan ambas al cuadrado, no es una parábola. Verificar que
$d(\text{vértice,foco})=d(\text{vértice,directriz})$: si no coinciden, hay error de
cuenta.
\end{atencion}

\subsection{Hipérbola}

\begin{teoria}
$\dfrac{x^2}{a^2}-\dfrac{y^2}{b^2}=1$ (horizontal) o $\dfrac{y^2}{a^2}-\dfrac{x^2}{b^2}=1$
(vertical), con $c^2=a^2+b^2$ (aquí $c$ es la hipotenusa, al revés que en la elipse).
Asíntotas: reemplazar el ``$1$'' por ``$0$'' y factorizar la diferencia de cuadrados.
\end{teoria}

\begin{proc}
\begin{enumerate}
    \item Llevar a forma canónica.
    \item El término \textbf{positivo} define el eje transversal (vértices y focos).
    $a^2$ = denominador del término \textbf{positivo} (no importa si es mayor o menor
    que $b^2$).
    \item $c=\sqrt{a^2+b^2}$. Asíntotas: poner ``0'', factorizar como diferencia de
    cuadrados e igualar cada factor a 0.
\end{enumerate}
\end{proc}

\begin{atencion}
\begin{itemize}
    \item A diferencia de la elipse, $a^2$ \textbf{no} es necesariamente el mayor
    denominador: es siempre el del término \textbf{positivo}. Aplicar la regla de la
    elipse acá es un error clásico.
    \item Asíntotas de la hipérbola horizontal: $y=\pm\frac ba x$. De la vertical:
    $y=\pm\frac ab x$ (¡el cociente se invierte!).
    \item Si el centro no es el origen, las asíntotas pasan por el centro $(h,k)$.
\end{itemize}
\end{atencion}

\subsection{Superficies cuádricas}

\begin{teoria}[Reconocimiento por signos, forma $Ax^2+By^2+Cz^2=R$ o variable lineal]
\begin{itemize}
    \item 3 signos iguales, $R\neq0$ del mismo signo: \textbf{elipsoide} (esfera si
    $A=B=C$).
    \item 2 positivos + 1 negativo, $=1$: \textbf{hiperboloide de una hoja}.
    \item 1 positivo + 2 negativos, $=1$: \textbf{hiperboloide de dos hojas}.
    \item Suma de 2 cuadrados = variable lineal: \textbf{paraboloide elíptico}.
    \item Diferencia de 2 cuadrados = variable lineal: \textbf{paraboloide hiperbólico}
    (``silla de montar'').
    \item Suma/diferencia de cuadrados $=0$ (dos positivos, uno negativo):
    \textbf{cono elíptico}.
    \item Falta una variable en la ecuación: \textbf{cilindro} (elíptico si mismo signo
    entre las dos presentes, hiperbólico si signos opuestos).
\end{itemize}
\end{teoria}

\begin{proc}
\begin{enumerate}
    \item ¿Hay una variable \textbf{lineal} (sin elevar al cuadrado) del lado derecho?
    Familia ``sin centro'' (paraboloides).
    \item Si todas están al cuadrado: contar signos positivos vs.\ negativos y comparar
    con la tabla.
    \item Ayuda visual: hacer secciones fijando $z=k$ (o $x=k$, $y=k$) y ver qué curva 2D
    resulta (elipse, hipérbola, punto, recta, vacío) para confirmar.
\end{enumerate}
\end{proc}

\begin{atencion}
\begin{itemize}
    \item Distinguir hiperboloide de \textbf{una} hoja (1 signo negativo) de \textbf{dos}
    hojas (2 signos negativos) es el error más común: en $z=0$, el de una hoja da una
    curva real (pasa por el ``cuello''); el de dos hojas no tiene puntos ahí.
    \item Si falta una variable por completo, es un \textbf{cilindro} que se extiende
    infinitamente en esa dirección — no es un error, es la definición.
    \item Ante la duda entre cono y paraboloide hiperbólico, verificar con 2 o 3
    secciones antes de decidir.
\end{itemize}
\end{atencion}

\end{document}
